\documentclass[12pt]{article}

\usepackage[T1]{fontenc}
\usepackage[utf8]{inputenc}
\usepackage[margin=1in]{geometry}
\usepackage{setspace}
\usepackage{parskip}
\usepackage{titlesec}
\usepackage[hidelinks]{hyperref}

\titleformat{\section}{\normalfont\bfseries\large}{}{0pt}{}
\titlespacing*{\section}{0pt}{1.4em}{0.6em}

\setstretch{1.15}

\begin{document}

\begin{center}
	{\Large\bfseries Primary Source Worksheet \#2}\\[0.6em]
	RELI 252 \quad\textbullet\quad Dr.\ Clarke\\[0.3em]
	Greyson Knippel
\end{center}

\vspace{1em}
\hrule
\vspace{1.5em}

\noindent\textbf{Primary Source Title:} Interview with Sister Kelly, from \textit{Unwritten History of Slavery: Autobiographical Accounts of Negro Ex-slaves}

\vspace{0.5em}

\noindent\textbf{Primary Source Author:} Sister Kelly (interviewed by Ophelia Settle Egypt, 1929)

\section*{What is the author's intent with this piece?}

Sister Kelly wants the person listening to her to actually understand what slavery was like, and she wants them to stop believing the softened/simplified version of it. She is not trying to entertain anyone. She says early on that she will never forget what happened ``this side of the grave,'' which tells you she sees remembering as a responsibility, not just a story she is telling.

A big part of what she wants to convince her listener of is that there was no such thing as a good enslaver. She brings up Mary Booker, the widow who owned her, and calls her ``the best white woman that ever broke bread.'' Then she immediately takes it back by saying ``that wasn't much, 'cause they all hated the po' ******.'' She proves her point with a detail instead of an argument. When Booker was dying, she only felt bad about leaving her children behind and ``didn't say a word 'bout her po' slaves what had done wuked theyselves to death for her.''

She also wants to convert the person she is talking to. In the middle of describing her life she stops and asks the interviewer if she knows the Lord, then tells her, ``you better know him, that's what'll help you.'' She is not just answering questions. She is preaching a little.

Finally, she wants her listener to know this is not ancient history. She says people ``stealed just like they stealed now, times ain't changed much.'' So she's speaking in 1929 and connecting slavery to what is still happening around her.

\section*{Compare and contrast with a previously assigned primary source}

Sister Kelly's interview and Francis Le Jau's letters describe the same religion from completely opposite sides. Le Jau was a white Anglican missionary in South Carolina writing reports back to his bosses in England in the early 1700s. Meanwhile, Sister Kelly was an enslaved girl in Tennessee remembering her life over a hundred years later. Both are talking about how enslaved people became Christian, but they describe two extremely different things.

Le Jau's version of Christianity is supervised. He only baptizes people with their master's permission, he makes them promise baptism will not make them want freedom, and he decides they probably should not learn to read after one man's Bible reading spread as prophecy. Everything happens where he can watch or control it. Sister Kelly's version happens in secret. She and the others turned a wash pot upside down in the middle of the floor to trap the sound so the white people could not hear them singing and praying. They were hiding the exact thing Le Jau wanted to organize.

And that difference explains a lot about how African American Christianity actually developed. Le Jau wanted religion to make enslaved people obedient, and the SPG advertised it that way. But Sister Kelly's practices and conversions had nothing to do with obeying anyone. She got religion at twelve as something that happened to her, and she says it gave her a faith that ``will stand for all time.'' It made her more confident, not more controllable.

The other big difference is who they are talking to. Le Jau is writing to impress white superiors. Sister Kelly is talking to a Black researcher from Fisk. She has no reason to make slavery sound better than it was, and she does not.

\end{document}